\section{The literal source and the imported foundations}\label{sec:setting}

Fix a computable literal parsing policy and a primitive recursive
control $f$. Write
\[
 F_0(n)=2^n,\qquad F_{j+1}(n)=F_j^{\,n+1}(n),\qquad
 G(n)=F_5(n+100).
\]
We assume $f\ge G$, $f(n+1)\ge2^{2^{f(n)}}$, and an eventual
upper bound by a fixed finite hierarchy level. The standard $G$, the
controls of \cite{groundcomparison}, and the additive-one controls
constructed below satisfy these assumptions. They are sufficient
conditions within the control freedom of \cite[Section~2.1]{ww}.

Ordinals below $\varepsilon_0$ are represented by hereditary Cantor
normal forms, with additive norm
\[
 N(0)=0,\qquad
 N\left(\sum_i\omega^{\alpha_i}c_i\right)
       =\sum_i c_i(1+N(\alpha_i)).
\]
Natural sum and product are $\oplus$ and $\otimes$.
The collapse is
\[
 \psi_f(\alpha)=\max\bigl(\{0\}\cup
 \{\psi_f(\beta)+1:\beta<\alpha,
                       N(\beta)\le f(N(\alpha))\}\bigr).
\]
All finite scalar additions are ordinary natural-number additions.
Finite tails in displayed Cantor forms use ordinary ordinal addition.

\subsection{Types, trees and the canonical clauses}

Simple types are generated by $o$ and $\sigma\to\tau$, with
$\lv(o)=0$ and
$\lv(\sigma\to\tau)=\max\{\lv(\sigma)+1,\lv(\tau)\}$.
The original constants have types
\[
 0:o,\quad S:o\to o,\quad D^\tau:o\to\tau\to\tau\to\tau,
 \quad R^\tau:o\to(o\to\tau\to\tau)\to\tau\to\tau.
\]
In particular $R^\tau0AB$ reduces to $B$, and
$R^\tau(Sn)AB$ reduces to $An(R^\tau nAB)$. Source terms use
ordinary typed application and abstraction. They are never normalized
before assignment.

An ordinal coordinate is a literal tree with leaves $1,\omega$ and
components $x_i$ of variable vectors, and binary constructors
\[
 \operatorname{Add}(a,b),\quad\operatorname{ExpProd}(a,b),
 \quad\operatorname{Psi}(a,b),
\]
interpreted by $a\oplus b$, $2^a\otimes b$ and
$\psi_f(\omega\otimes a\oplus b)$. The latter two are
primitive binary constructors. An actual graded coordinate contains
no Zero node. A vector of level $m$ has $m+1$ actual coordinates;
out-of-range coordinates are absent, and only in comparisons are they
interpreted as zero. Vector addition retains a single actual operand
unchanged and joins two actual operands by one Add node.

The policy chooses an ordered full binary addition tree at each
nonempty finite-sum site and arity. Every operand is used once, in
order and unchanged; a positive repeat uses exactly its designated
number of copies. Choices do not depend on names, values, inner trees
or free-variable status. Existing subtrees are not reassociated.
The sites include positive numerals, repeats, active three-operand
partials, positive abstraction sums and the optional three-operand
box tail. Binary clauses stay binary. In an exponential partial the
two-fold head repeat remains one nested operand of the outer sum.
Appendix~\ref{app:policy} shows why this convention affects values.

For $a,b$ of levels $m>n$, the box $c=a\square b$ is
\[
 c_i=a_i\ (n<i\le m),\qquad
 c_i=\operatorname{ExpProd}(c_{i+1},a_i+b_i)\ (1\le i\le n),
 \qquad c_0=\operatorname{Psi}(c_1,a_0+b_0+n).
\]
The finite $n$ contribution is absent when $n=0$ and is a positive
numeral otherwise. All finite macros follow the policy.

For a tracked vector $\mathbf x$, put $n=\lv(X)+1$. Whole-freeness
means absence of every component of $\mathbf x$. The size
$s_X(h)$ is one if $h$ is whole-free or is a tracked component.
Otherwise it is $s_X(a)+s_X(b)+1$ at Add and Psi, and
$2s_X(a)+s_X(b)+1$ at ExpProd. The graded partial
$\partial_i^Xh$ has level $n$. Its whole-free case, tested first,
has diagonal $h+1$ if $i\le n$ and one elsewhere. The active clauses are
\begin{align*}
 \partial_i^Xx_i&=\mathbf1,\\
 \partial_i^X(a+b)&=\partial_i^Xa+\partial_i^Xb+\mathbf1,\\
 \partial_i^X\operatorname{ExpProd}(a,b)
      &=2\partial_{i+1}^Xa+\partial_i^Xb+\mathbf1\quad(i>0),\\
 (\partial_0^X\operatorname{Psi}(a,b))_0
      &=\operatorname{Psi}(a\{\mathbf x:=\mathbf1\},(\partial_0^Xb)_0),\\
 (\partial_0^X\operatorname{Psi}(a,b))_j
      &=(\partial_1^Xa)_j+(\partial_0^Xb)_j\quad(1\le j\le n).
\end{align*}
For $h$ of level $m$, let $d_X(h)=2^n\sum_{i=0}^m s_X(h_i)$.
The abstraction operator has level $\max(n,m)$ and
\[
 (\delta^Xh)_0=d_X(h)(\partial_0^Xh_0)_0,\quad
 (\delta^Xh)_j=\sum_{i=0}^m(\partial_i^Xh_i)_j\ (0<j\le n),
 \quad(\delta^Xh)_j=h_j\ (j>n).
\]
The multiplier is computed from the fixed literal input trees.

The canonical assignment is the zero-history instance of
\cite[Definition~3.1]{ww}:
\begin{align*}
 \can{X:\sigma}&=(x_0,\ldots,x_{\lv(\sigma)}),&
 \can0&=(1),&\can S&=(1,1),\\
 \can{D^\tau}&=(1,\ldots,1),&&&
 \lv(\can{D^\tau})&=\lv(\tau)+1,\\
 \can{R^\tau}&=(2,1,\ldots,1,\omega),&&&
 \lv(\can{R^\tau})&=\lv(\tau)+2,\\
 \can{AB}&=(\can A\square\can B)\upharpoonright_{\lv(\tau)}
                  &&(A:\sigma\to\tau, B:\sigma),\\
 \can{\lambda X.A}&=\delta^X\can A+
                           \can A\{\mathbf x:=\mathbf1\}.
\end{align*}
All introduced bound names below are fresh. A natural input $r$ is
encoded by the positive ground coordinate $r+1$; source $0$ therefore
encodes arithmetic zero. Unless a common unit observation is stated,
all free source names in comparison queries are ground. Weak comparison
is universal coordinatewise $\le$; strict comparison requires universal
$<$ in coordinate zero and universal $\le$ in the positive coordinates.

\subsection{Exact foundational interface}

We use the following proved interface from the repaired canonical
assignment theorem and closed evaluator in \cite{groundcomparison}.
Every actual assignment is a finite effectively constructed graded
vector of its type level, with positive coordinates and bounded norm
$N(h_i)\le h_0$ under admissible valuations. The partials have their
graded domains and $\delta$ retains the multiplier printed above.
Its scalar estimate is the weak inequality
$h\{\mathbf x:=\mathbf1\}\le(\partial_0^Xh)_0$.
No strict inequality at the tracked scalar variable is assumed.

Generated positive coordinates contain only positive-index variable
components, $1,\omega$, Add and graded ExpProd. They contain neither
scalar components nor Psi. They are therefore independent of every
ground input, and become closed after all free functional names receive
the common unit observation. Such closed ordinals and all closed scalar
values are computable.

For a nonzero limit $\lambda$ define
\[
 J_\lambda(t)=\psi_f(\lambda+t-N(\lambda)),\qquad t\ge N(\lambda).
\]
The exact evaluator replaces a nonfinite state $\alpha$ by
$M(\limpart(\alpha),f(N(\alpha)))$, where
$M(\eta,b)=\max\{\gamma<\eta:N(\gamma)\le b\}$.
Every new norm is exactly $f$ of the preceding norm. If $D_\lambda(t)$
counts these nonfinite events and $Z_\lambda(t)$ is the first finite
state, then
\[
 J_\lambda(t)=Z_\lambda(t)+D_\lambda(t),\qquad
 Z_\lambda(t)=f^{D_\lambda(t)}(t).
\]
In particular $\psi_f(\omega k+n)=f^k(2k+n)+k$.
At a fixed limit head the collapse is strictly increasing in its
finite tail. These statements have their complete proofs in the cited
companion; no comparison decision theorem is needed to construct the
sources in this paper.

The policy also makes substitution covariance literal. Substitution for
another name by a tracked-name-free vector preserves whole-freeness,
the active constructor pattern and $s_X$. Induction through the displayed
partials proves their covariance and then that of $\delta$.
Consequently capture-free replacement of a ground source variable by
literal zero commutes with canonical assignment. This is the instance
of \cite[Lemma~2.20]{ww} used by the observers below.
